% OUP official authoring-template LaTeX package for Briefings in Bioinformatics.
% Generated from submission_package/01_manuscript/main_manuscript.md (v0.6, real-data recompute pass).
% OUP guidance: use the OUP LaTeX template with the 'numbered' option for Vancouver-style citations.
\documentclass[modern,large,numbered,nocrop]{oup-authoring-template}

\usepackage{booktabs}
\usepackage{tabularx}
\usepackage{url}
\usepackage{hyperref}
\usepackage{float}
\graphicspath{{figures/}}
\setcitestyle{numbers,square,comma}
\makeatletter
\@numbibtrue
\NAT@numberstrue
\makeatother

% Article type (select in ScholarOne at submission; also declared in the cover letter): Problem Solving Protocol
\journaltitle{Briefings in Bioinformatics}
\pubyear{2026}
\copyrightyear{2026}
\def\societylogo{}% no society logo image supplied; suppresses the template's default PNAS-Nexus-Logo-CMYK asset
\title[Transcriptional reach of the NK serine--SM--topology axis]{Transcriptional Reach of the Serine--Sphingomyelin--Membrane-Topology Axis in NK-Cell Immune Evasion: A Single-Cell Heterogeneous Graph Framework from Liver to Gastric Cancer}

\author[1,2]{Guohao Lyu}
\author[1,2]{Yingchun Xia}
\author[1,2]{Huichao Liu}
\author[1,2]{Xiaolei Zhu}
\author[1,2]{Shuai Yang}
\author[3,4,*]{Ailian Zhou}
\author[1,2,*]{Lichuan Gu}
\address[1]{School of Artificial Intelligence, Anhui Agricultural University, Hefei 230036, China}
\address[2]{Anhui Province Key Laboratory of Intelligent Agricultural Technology and Equipment, Anhui Agricultural University, Hefei 230036, China}
\address[3]{Agricultural Information Institute, Chinese Academy of Agricultural Sciences, Beijing 100081, China}
\address[4]{Key Laboratory of Agricultural Blockchain Application, Ministry of Agriculture and Rural Affairs, Beijing 100081, China}
\corresp[*]{Correspondence should be addressed to Lichuan Gu. Email: glc@ahau.edu.cn}
\corresp[*]{Correspondence should be addressed to Ailian Zhou. Email: zhouailian@caas.cn}
% ORCID iDs: by author decision, not included in the manuscript text; each
% author will enter their iD directly in Oxford ScholarOne at submission.

\abstract{\subheading{Motivation} Zheng et al. (\textit{Nat Immunol} 2023) showed that tumors evade natural killer (NK) cell cytotoxicity by depleting serine-derived sphingomyelin (SM) in NK membranes, collapsing membrane protrusions and abolishing lytic synapse formation -- a mechanism proven with single-cell mass spectrometry and super-resolution imaging in a small liver-cancer cohort, methods that cannot scale to cohort-level analysis or be tested in other cancers. \subheading{Results} We present GC-NKGraph-Atlas, which asks how much of this serine$\to$sphingomyelin$\to$membrane-topology$\to$cytotoxicity axis is transcriptionally recoverable from public tumor transcriptomes. A \textit{mechanism-card} formalism encodes the wet-lab mechanism as a machine-readable recipe driving cell-type-attributed proxy scoring, a mechanism-grounded heterogeneous gene graph, and a graph neural network (GNN) that predicts NK immune states. Across a two-arm design (liver positive control, gastric extension) we obtain a three-layer scoping result. \textbf{First}, the \textit{effector arm} (protrusion-machinery$\leftrightarrow$cytotoxicity) is detectable in bulk transcriptomes (TCGA-LIHC r=0.55; gastric replications GSE62254 r=0.42, GSE84437 r=0.62), but roughly half of this coupling reflects NK-cell-abundance covariation -- after adjusting for an NK-fraction proxy it survives in two of three bulk cohorts, not the third, and the single-cell version of the same coupling does not survive residualization against library size/latent structure or a random-gene-module permutation baseline. We therefore anchor the effector-arm claim on the bulk result, treated as partially recoverable and substantially infiltration-driven. \textbf{Second}, transcription does not proxy the physical topology phenotype: intratumoral NK cells show \textit{higher}, not lower, protrusion-machinery transcripts ($\Delta$=+0.14) despite reduced cytotoxic output ($\Delta$=$-$0.14) -- opposite to the physical membrane collapse. \textbf{Third}, the upstream metabolic arm (SM-balance$\to$protrusion) is not transcriptionally detectable at either resolution (bulk r=$-$0.02; single-cell corrected r=0.03, P=0.20), consistent with metabolite-level regulation. A second, independently authored mechanism card (TGF$\beta$$\to$SMAD$\to$NK exclusion) reproduces the same boundary end-to-end, generalizing the finding across two distinct evasion mechanisms. The framework prioritizes \textbf{37 putative tumor-intrinsic candidate targets} (led by PHGDH, SGMS2, PSAT1, PSPH, SMPD3/1), each with a recommended validation assay, presented as a hypothesis-generation shortlist rather than a tumor-exclusive target set. \subheading{Availability} All code, configuration, and synthetic test data are available at \url{https://github.com/nblvguohao/GC-NKGraph-Atlas}. The mechanism-card template enables application to additional published mechanisms without modifying the pipeline core.}
\keywords{NK cell; immune evasion; heterogeneous graph neural network; single-cell transcriptomics; sphingomyelin; gastric cancer}

\begin{document}
\maketitle
\enlargethispage{3\baselineskip}% avoid abstract text colliding with the title-page footer

\section*{Key Points}
\begin{itemize}
\item A \textbf{mechanism-card} formalism converts published wet-lab immune-evasion mechanisms into scalable transcriptome analyses; demonstrated end-to-end on two independent cards (serine--SM; TGF$\beta$$\to$SMAD) that reveal a consistent transcriptional-reach boundary.
\item A \textbf{two-arm design} (liver positive control + gastric extension) shows the effector arm (protrusion$\to$cytotoxicity) is recoverable in bulk transcriptomes but \textbf{$\sim$50\% of the coupling is NK-abundance covariation}, and the single-cell coupling does not survive technical-confound controls; the metabolic arm is not detectable after correction (p=0.20).
\item \textbf{The reach boundary generalizes:} a second, independently authored card reproduces the same pattern -- generic NK-effector co-variation recovers, but mechanism-specific causal predictions fail in bulk.
\item \textbf{Transcription does not proxy the physical topology phenotype:} intratumoral NK cells show higher protrusion-machinery transcripts despite collapsed physical protrusions -- a fundamental disconnect defining the boundary of transcriptome-based reconstruction.
\item The graph is used as a \textbf{probe, not a predictor}: ablating the \url{metabolic_crosstalk} edge abolishes the corresponding embedding coupling, converging with the null correlation result; the framework outputs a de-circularized 37-gene candidate list with recommended validation assays.
\end{itemize}

\section{Introduction}

Natural killer (NK) cells kill tumor cells without prior antigen sensitization, making them attractive cancer-immunotherapy effectors [1--5], but tumors evade NK killing through metabolic competition, inhibitory-ligand upregulation, immunosuppressive signaling, and physical exclusion [6--11]. Zheng et al. [12] elucidated a mechanism in which tumors dysregulate serine metabolism [13,14], depleting NK-membrane sphingomyelin (SM) [15,16], collapsing membrane protrusions, and abolishing lytic synapse formation -- rescuable by SM-catabolism inhibition, synergistic with Tim3 blockade, and proposed to extend beyond liver cancer [17]. However, this mechanism was proven using super-resolution imaging and single-cell mass spectrometry in a small cohort with no deposited transcriptomic data -- techniques that do not scale to cohort-level testing.

Three gaps motivate this work: (i) \textbf{scalability} -- testing the axis across hundreds of samples requires a transcriptome-based proxy; (ii) \textbf{transferability} -- whether the axis operates in gastric cancer, not on the mechanism's published extension list, is unknown; (iii) \textbf{reusability} -- each mechanistic discovery currently requires a bespoke computational follow-up.

We introduce \textbf{GC-NKGraph-Atlas}, built around a reusable \textbf{mechanism-card} abstraction: a machine-readable YAML specification encoding one published mechanism's molecular chain, gene modules, expected directions, physical ground-truth targets, and validation logic. The card driving this study (\url{zheng_nk_sm_topology}) drives (a) cell-type-attributed transcriptional proxies from scRNA-seq, (b) a tumor--NK heterogeneous graph with a \url{metabolic_crosstalk} edge grounded in the specific biology, (c) a graph neural network learning gene embeddings to predict NK immune states, and (d) multi-evidence candidate-target ranking. A two-arm design tests \textbf{Arm A} (liver/HCC positive control, the system where the axis was proven) and \textbf{Arm B} (gastric-cancer extension).

Our framework draws on immune-state/cell-type inference tools (CIBERSORTx [18], quanTIseq [19], Scissor [20]), cell--cell-communication and metabolic-flux inference (CellChat [22], CellPhoneDB [23], NicheNet [24], scFEA [25]), trajectory/dysfunction modeling (CytoTRACE [26]; CREM regulation [27]), and graph learning in cancer genomics [28--32], including genome-scale driver-gene-discovery fusion models such as TREE [49] and GRAFT [50]. None of these operationalizes a single published mechanism as a machine-readable card with pre-registered claim boundaries between transcriptional proxy and physical phenotype, nor asks -- rather than assumes -- how far a specific mechanism's transcriptional signal reaches; TREE and GRAFT in particular target genome-scale, cancer-agnostic driver discovery from generic networks, a different scale, label, and objective from this study's $\sim$100-gene, mechanism-anchored, reach-mapping design (\S{}4).

\section{Methods}

\textbf{Study design.} The pipeline runs 14 phases across five stages: data acquisition (TCGA-STAD/LIHC, GEO gastric cohorts), scRNA integration and NK-state scoring, heterogeneous graph construction, GNN model training against six tabular baselines, and target prioritization (\url{src/pipeline.py}, checkpoint-based, with a synthetic-data test mode).

\begin{figure*}[htbp]
\centering
\includegraphics[width=\textwidth]{figures/fig0_workflow.pdf}
\caption{Workflow overview of the GC-NKGraph-Atlas framework, from data acquisition through target prioritization, showing the mechanism-card abstraction and the two-arm study design (Arm A: liver positive control; Arm B: gastric cancer extension). \textit{Alt text: Workflow overview of the GC-NKGraph-Atlas framework showing the computational pipeline stages, the mechanism-card abstraction, and the two-arm study design.}}\label{fig:workflow}
\end{figure*}

\textbf{Mechanism-card abstraction.} A mechanism card is a YAML specification (schema: origin, biology, transcriptional\_proxy, physical\_ground\_truth [gated], graph\_integration, validation, therapeutic\_hook; full schema in \url{configs/mechanism_cards/mechanism_card.template.yaml}) that separates what is computable from transcriptomes from what requires physical measurement; mock data are excluded from all scientific claims.

\textbf{SST-axis transcriptional proxy.} The serine--SM--topology axis is operationalized as seven gene modules ($\sim$100 genes: tumor serine capacity, NK SM synthesis/catabolism, NK de-novo sphingolipid, NK protrusion machinery, NK synapse/cytotoxicity outcome, checkpoint link) yielding derived scores \url{nk_sm_balance}, \url{nk_topology_permissive}, and \url{sst_axis_score}. Claims are strictly limited to ``transcriptional program permissive-of/associated-with'' the topology phenotype; physical-topology or metabolite-content prediction is prohibited.

\textbf{scRNA-seq analysis.} GSE246662 (nine samples, healthy liver/gastric cancer/liver metastasis, 166,829 cells) was library-size normalized, scVI-integrated (30 latent dims), and Leiden-clustered (SCANPY [33--37]); NK cells (n=8,310) were separated from T/monocyte/B lineages by marker threshold. Four NK states (hot-cytotoxic, hot-dysfunctional, cold/excluded, intermediate) were defined from cytotoxicity/dysfunction/exclusion gene-set z-scores [39,40] and projected onto bulk samples by scRNA-anchored scoring. No hard per-cell QC threshold was applied (Limitations); technical variance is addressed via count-depth residualization (\S{}3.2).

\textbf{Heterogeneous graph.} Nodes are the $\sim$100 SST-axis/candidate genes; edges combine standard priors (STRING PPI, CellChatDB ligand--receptor) with two mechanism-grounded types (\url{metabolic_crosstalk}: tumor-serine gene to NK-topology gene, weight 0.5, fixed not fit to data; \url{sm_topology_axis}) and scRNA co-expression. The \url{metabolic_crosstalk} edge is diagnostic, not predictive: it lets us ask whether hard-coding the mechanism's upstream coupling is corroborated by the data (\S{}3.7).

\textbf{Model.} Gene embeddings are learned by truncated-SVD spectral decomposition of the combined adjacency matrix (all reported results, including the ablation, use this encoder; a learnable heterogeneous graph transformer [32, 41] was not used -- \S{}4.4). A classifier MLP on [expression, expression$\cdot$embedding] predicts NK state (5-fold stratified CV; Adam, hyperparameters selected by a 100-trial Bayesian search). Six tabular baselines (XGBoost [42], LightGBM [43], Random Forest, ElasticNet, SVM, MLP) and two domain-knowledge baselines (NK-marker signature; SST-module signature) are evaluated on identical splits. A six-view label-masked multi-view audit (generic/LR/\url{metabolic_crosstalk}/\url{sm_topology_axis}/GO/MSigDB views [49,50], spectral SVD, global-softmax fusion) tests whether learned view-fusion adds predictive value once every NK-state-label-defining gene is masked from inputs, using ten fixed seeds (STAD 80:20 train/early-stop, LIHC held out).

\textbf{Target prioritization.} Candidate genes are ranked by a five-dimension composite (tumor specificity 0.30, NK-dysfunction correlation 0.20, SST-axis membership 0.30, axis-core membership 0.10, literature support 0.10); tumor-intrinsic candidates are required to be up in malignant cells, separating them from NK-side axis-readout genes.

\textbf{Pre-registered hypotheses and pseudoreplication correction.}

\begin{table*}[htbp]
\caption{Pre-registered hypotheses.}
\small
\begin{tabularx}{\textwidth}{@{}lXX@{}}
\toprule
ID & Hypothesis & Expected direction \\
\midrule
H1 & tumor\_serine\_capacity $\perp$ nk\_sm\_balance & Direction calibrated, reported \\
H2 & nk\_sm\_balance (+) nk\_protrusion\_machinery & Positive \\
H3 & nk\_protrusion\_machinery (+) cytotoxicity anchors & Positive \\
H4 & nk\_topology\_permissive ($-$) HAVCR2 / dysfunction & Negative \\
H5 & intratumoral NK < peritumoral/peripheral NK in sm\_balance \& protrusion machinery (scRNA) & Negative (tumor < normal) \\
\bottomrule
\end{tabularx}
\end{table*}

Five hypotheses were pre-registered for Arm A with a full-recovery criterion requiring H2--H5 to pass. Single-cell NK correlations (8,310 cells, 9 samples) are corrected for pseudoreplication by per-sample Pearson r, Fisher-z transform, and DerSimonian-Laird random-effects meta-analysis across samples; corrected values are reported throughout (naive per-cell values retained in \url{results/tables/sst_axis_pseudoreplication_corrected.tsv} for transparency). Full protocol detail (QC thresholds, scVI/Leiden parameters, hyperparameter search, and the multi-view audit design) is in Supplementary Methods S.M.1--S.M.5.

\section{Results}

\subsection{Dataset summary}

Four bulk cohorts and one multi-tissue scRNA dataset were processed (Table~\ref{tab:datasets}): TCGA-LIHC (n=423, liver, positive control), TCGA-STAD (n=450, gastric, primary training [44]), GSE62254/ACRG (n=300 [45]) and GSE84437 (n=483, gastric, external validation; both required probe-to-symbol remapping), and GSE246662 (166,829 cells, 8,310 NK). Bulk NK-state labels were assigned by scRNA-anchored scoring.

\begin{table*}[htbp]
\caption{Dataset characteristics.}\label{tab:datasets}
\footnotesize
\begin{tabularx}{\textwidth}{@{}lXXXX@{}}
\toprule
Dataset & Cancer & Type & Samples / cells & Role \\
\midrule
TCGA-LIHC & Liver & Bulk RNA-seq & 423 & Positive control (Arm A) \\
TCGA-STAD & Gastric & Bulk RNA-seq & 450 & Train primary (Arm B) \\
GSE62254 (ACRG) & Gastric & Bulk microarray & 300 & External validation \\
GSE84437 & Gastric & Bulk microarray & 483 & External validation \\
GSE246662 & Liver/gastric/metastasis & scRNA-seq & 166,829 cells (8,310 NK) & NK atlas + axis test \\
\bottomrule
\end{tabularx}
\end{table*}

\subsection{Arm A -- the axis is partially recovered: effector arm yes, topology arm no}

Hypotheses H1--H5 were tested at bulk and single-cell resolution (Table~\ref{tab:hyp-outcomes}; Figure~\ref{fig:armA}). \textbf{H1} (serine-capacity$\leftrightarrow$SM-balance) is null at both resolutions (bulk r=$-$0.02, P=0.74; single-cell r=+0.01, P=0.27), consistent with serine$\to$SM crosstalk operating at the metabolite rather than transcript level. \textbf{H3} (protrusion$\leftrightarrow$cytotoxicity, the effector arm) recovers in bulk TCGA-LIHC (r=0.55, P=5$\times$$10^{-35}$), but two controls qualify this. First, roughly half of the bulk coupling is NK-cell-abundance covariation: partialling out a clean NK-lineage-marker proxy attenuates it to a partial r=0.25 in TCGA-LIHC, holding in two of three bulk cohorts but abolished in the third (GSE84437, partial r=$-$0.07, n.s.). Second, the same coupling at single-cell resolution (corrected meta-analytic r=0.31 across 9 samples) does not survive further technical-confound controls: it collapses to r=0.09 after residualizing against library size and the real scVI latent space, and does not exceed a random-gene-module permutation null (empirical P=0.97). We therefore anchor the effector-arm claim on the bulk result and its gastric replications (\S{}3.3), and report the effector arm as \textit{partially recoverable and substantially infiltration-driven}. \textbf{H2} (SM-balance$\to$protrusion, the metabolic arm) is undetectable in bulk (r=$-$0.02, P=0.72) and not significant after pseudoreplication correction at single-cell resolution (corrected r=0.03, P=0.20). \textbf{H4/H5} (physical topology) show the most consequential result: intratumoral NK cells have \textit{higher}, not lower, protrusion-machinery transcripts than normal-tissue NK ($\Delta$=+0.14, P=3$\times$$10^{-91}$) despite correctly reduced cytotoxic output ($\Delta$=$-$0.14, P=6$\times$$10^{-52}$) -- opposite to the physical membrane collapse, externally corroborated for the cytotoxicity direction [46]. This is the natural boundary of transcriptome-based reconstruction for a mechanism whose phenotype is physical and whose causal link operates at the metabolite level, precisely why the anchor study required single-cell mass spectrometry and super-resolution imaging. Full derivation for H2--H5 is in Supplementary Results S3.2.

\begin{table*}[htbp]
\caption{Pre-registered hypothesis outcomes (multi-resolution). \textdagger{} Single-cell NK r/P are meta-analytic values after pseudoreplication correction (Methods \S{}2.9). \textdaggerdbl{} The H3 single-cell coupling additionally fails count-depth/scVI-latent residualization and a random-module permutation test (empirical P=0.97); full battery in Supplementary Results S3.2. The effector-arm claim rests on the bulk result (this table) and its bulk gastric replications (\S{}3.3), not the single-cell H3 number.}\label{tab:hyp-outcomes}
\scriptsize
\resizebox{\textwidth}{!}{\begin{tabular}{@{}p{0.7cm}p{2.2cm}p{1.3cm}p{0.9cm}p{1.0cm}p{1.0cm}p{2.6cm}@{}}
\toprule
Hyp & Test & Resolution & r / $\Delta$ & $P$ & Expected & Outcome \\
\midrule
H1 & serine\_cap \textasciitilde{} sm\_balance & bulk & $-$0.016 & 0.74 & calibrated & reported (null) \\
H1 & serine\_cap \textasciitilde{} sm\_balance & single-cell NK & +0.012 & 0.27 & calibrated & reported (null) \\
H2 & sm\_balance \textasciitilde{} protrusion & bulk & $-$0.017 & 0.72 & + & not recovered \\
H2 & sm\_balance \textasciitilde{} protrusion & single-cell NK & +0.029 & 0.20 & + & not signif.\ after correction \\
H3 & protrusion \textasciitilde{} cytotoxicity & bulk & +0.551 & 5$\times$$10^{-35}$ & + & \textbf{recovered} \\
H3 & protrusion \textasciitilde{} cytotoxicity & single-cell NK & +0.313\textdagger\textdaggerdbl & 3.9$\times$$10^{-8}$\textdagger\textdaggerdbl & + & not distinguishable from technical background \\
H4 & topology \textasciitilde{} dysfunction & bulk & +0.311 & 7$\times$$10^{-11}$ & $-$ & not recovered \\
H4 & topology \textasciitilde{} HAVCR2 & single-cell NK & +0.036\textdagger & 9.1$\times$$10^{-4}$\textdagger & $-$ & not recovered (after corr.) \\
H5 & intratumoral<normal: cytotoxicity & single-cell NK & $-$0.141 & 6$\times$$10^{-52}$ & tumor< & \textbf{recovered} \\
H5 & intratumoral<normal: protrusion & single-cell NK & +0.142 & 3$\times$$10^{-91}$ & tumor< & not recovered \\
\bottomrule
\end{tabular}}
\end{table*}

\begin{figure*}[htbp]
\centering
\includegraphics[width=\textwidth]{figures/fig1_armA_positive_control.pdf}
\caption{(a) Confound-control logic for the effector-arm claim. (b) H3: protrusion-machinery vs cytotoxicity-output score, bulk TCGA-LIHC (r=0.55, p=5$\times$$10^{-35}$). (c) H2 and H3 at bulk vs single-cell resolution. (d) Intratumoral vs normal NK per-cell module z-scores. (e) Table~\ref{tab:hyp-outcomes} as a forest plot. \textit{Alt text: Five-panel summary of liver positive-control recovery.}}\label{fig:armA}
\end{figure*}

\subsection{Arm B -- Gastric cancer extension}

The effector coupling extends to gastric cancer: scRNA (r=0.35, P=7$\times$$10^{-30}$, n=1,017 gastric NK, sharing the same single-cell technical caveats as \S{}3.2) and, after platform-annotation remapping (GSE62254$\to$GPL570; GSE84437$\to$GPL6947), both external microarray cohorts (Table~\ref{tab:extval}). As in Arm A, this bulk replication is substantially NK-abundance-driven (\S{}3.2): after partialling out NK fraction it remains significant in GSE62254 but not GSE84437. The metabolic coupling (H2), null in liver, is weakly but significantly positive in both gastric bulk cohorts (r=0.18/0.11), consistent with a real but small effect more detectable where NK signal is stronger. Gastric-cancer NK show the lowest cytotoxicity-output score of the three tissues sampled, consistent with an evasion phenotype (Figure~\ref{fig:armB}).

\begin{table*}[htbp]
\caption{External validation of the axis (independent gastric cohorts).}\label{tab:extval}
\footnotesize
\begin{tabularx}{\textwidth}{@{}lXXXX@{}}
\toprule
Cohort & n & protrusion\textasciitilde{}cytotoxicity r (p) & sm\_balance\textasciitilde{}protrusion r (p) & NK markers \\
\midrule
GSE62254 (GPL570) & 300 & 0.42 (1.4$\times$$10^{-14}$) & 0.18 (1.3$\times$$10^{-3}$) & 6/7 \\
GSE84437 (GPL6947) & 483 & 0.62 (3.3$\times$$10^{-53}$) & 0.11 (0.02) & 7/7 \\
\bottomrule
\end{tabularx}
\end{table*}

\begin{figure*}[htbp]
\centering
\includegraphics[width=\textwidth]{figures/fig2_armB_extension.pdf}
\caption{(a) External-validation design. (b) Cross-tissue module comparison. (c) NK-immune-state distribution in TCGA-STAD (n=450). (d) Effector-coupling replication across cohorts. \textit{Alt text: Four-panel gastric-cancer extension figure.}}\label{fig:armB}
\end{figure*}

\subsection{NK immune state classification}

On TCGA-STAD (5-fold CV, NK-hot-cytotoxic vs rest), the graph-informed GNN (MCC 0.706, AUROC 0.950) is statistically indistinguishable from the top gradient-boosting baselines (LightGBM MCC 0.733, XGBoost 0.727; paired P>0.28) and from a no-graph SST-module-signature baseline (MCC 0.619, AUROC 0.904; P>0.10), while significantly exceeding a simple NK-marker-signature baseline and the linear/kernel/shallow-network baselines (Table~\ref{tab:classification}, Figure~\ref{fig:model}). The graph therefore does not provide a discriminative advantage over simpler alternatives; its value here is diagnostic (\S{}3.7), not predictive. Full baseline detail, including a CIBERSORTx/quanTIseq/Scissor comparison roadmap, is in Supplementary Results S3.4.

\begin{table*}[htbp]
\caption{NK-state classification (TCGA-STAD, 5-fold CV; mean over folds). GNN not significantly different from LightGBM/XGBoost (paired p$>$0.27) or the SST-module signature baseline (p$>$0.10); significantly above the NK-marker signature baseline and ElasticNet/SVM/MLP (p$<$0.05).}\label{tab:classification}
\scriptsize
\resizebox{\textwidth}{!}{\begin{tabular}{@{}p{1.0cm}p{1.65cm}p{1.65cm}p{1.65cm}p{1.65cm}p{1.65cm}p{1.65cm}@{}}
\toprule
Method & Accuracy & Balanced Acc. & Macro F1 & MCC & AUROC & AUPRC \\
\midrule
LightGBM & 0.880 & 0.863 & 0.865 & \textbf{0.733} & \textbf{0.960} & 0.933 \\
XGBoost & 0.878 & 0.859 & 0.862 & 0.727 & 0.959 & 0.928 \\
\textbf{GC-NKGraph-Atlas (GNN)} & 0.864 & 0.856 & 0.850 & 0.706 & 0.950 & 0.910 \\
RandomForest & 0.860 & 0.825 & 0.836 & 0.681 & 0.941 & 0.901 \\
ElasticNet & 0.856 & 0.822 & 0.832 & 0.671 & 0.927 & 0.884 \\
SST-module signature & 0.818 & 0.823 & 0.803 & 0.619 & 0.904 & 0.827 \\
NK-marker signature & 0.787 & 0.744 & 0.750 & 0.503 & 0.849 & 0.733 \\
SVM (RBF) & 0.773 & 0.672 & 0.683 & 0.489 & 0.900 & 0.839 \\
MLP (2-layer) & 0.747 & 0.669 & 0.641 & 0.376 & 0.807 & 0.730 \\
\bottomrule
\end{tabular}}
\end{table*}

\begin{figure*}[htbp]
\centering
\includegraphics[width=\textwidth]{figures/fig4_model_comparison.pdf}
\caption{MCC (a) and AUROC (b) for GC-NKGraph-Atlas versus tabular and signature baselines on identical seed-42 folds; error bars show cross-fold SD. \textit{Alt text: Bar charts of MCC and AUROC comparing the graph model with tabular and signature baselines.}}\label{fig:model}
\end{figure*}

\subsection{Candidate target prioritization (de-circularized)}

Separating tumor-intrinsic candidates (required up in malignant cells) from NK-side axis-readout genes yields 37 tumor-intrinsic candidates, headed by the druggable serine/SM enzymes PHGDH, SGMS2, PSAT1, PSPH, and SMPD3/1 (Table~\ref{tab:targets}). A trivial mechanism-card-membership-only baseline correlates moderately with the five-dimension scoring (Spearman rho=0.54): the added value is re-ordering by tumor specificity and surfacing 12 gastric-relevant non-SST genes (e.g., COL1A1/2, NECTIN2, ERBB2, FGFR2, MET). Two orthogonal filters -- NK-state differential expression (externally replicated in two gastric cohorts) and DepMap 26Q1 CRISPR essentiality in 35 gastric cell lines -- refine the list into evidence tiers (Table~\ref{tab:targets}): \textbf{SMPD3/SMPD1} pass both filters cleanly; \textbf{SGMS2} and \textbf{NT5E/CD73} pass on distinct evidence but carry caveats (discordant external DE; label self-referentiality, respectively); five serine-pathway genes and RAC1/ERBB2 are downgraded or excluded for wrong-direction DE or pan-essentiality. The list is de-circularized but not fully clean (46\% of candidates retain NK-side module annotation) and is a hypothesis-generation shortlist, not a validated tumor-exclusive set. Full evidence and external replication are in Supplementary Results S3.5.

\begin{table*}[htbp]
\caption{Top putative tumor-intrinsic candidate targets in gastric cancer, evidence-tiered. Tier 1: non-pan-essential plus NK-state DE or mechanistic privilege. Tier 2: single-dimension support. Tier 3: downgraded (NK-state signal absent/opposite). Tier X: excluded (pan/moderately essential). \S{} SGMS2's DE signal is discordant in external replication. \textdaggerdbl{} NT5E's DE result is partly self-referential (Supplementary Results S3.5).}\label{tab:targets}
\scriptsize
\resizebox{\textwidth}{!}{\begin{tabular}{@{}p{0.9cm}p{1.6cm}p{2.1cm}p{1.4cm}p{2.3cm}p{2.5cm}p{2.0cm}@{}}
\toprule
Tier & Gene & Category & tumor\_spec (log2FC) & NK-state DE & DepMap CERES (real, 35 gastric lines) & Druggability \\
\midrule
1\S{} & SGMS2 & SM synthesis & +0.038 & \textbf{UP in dysf (FDR=0.007)}; discordant externally & weakly essential ($-$0.02) & Preclinical \\
1\textdaggerdbl & NT5E & adenosine (CD73) & +0.013 & \textbf{UP in dysf (FDR=0.014)} & non-essential (+0.005) & Phase 1/2 clinical \\
1 & SMPD3 & SM catabolism & +0.032 & no signal & weakly essential ($-$0.13) & Preclinical (Zheng 2023) \\
1 & SMPD1 & SM catabolism & +0.001 & no signal & weakly essential ($-$0.16) & Preclinical (Zheng 2023) \\
2 & PHGDH & serine metabolism & +0.059 & no signal & weakly essential ($-$0.04) & Phase 1/2 (PHGDH inhibitor) \\
2 & MICA & stress ligand & +0.006 & no signal & non-essential (+0.05) & Preclinical \\
3 & PSAT1 & serine metabolism & +0.016 & no signal & moderately essential ($-$0.27) & -- downgraded \\
3 & PSPH & serine metabolism & +0.010 & \textbf{DOWN in dysf (p=0.03)} & weakly essential ($-$0.16) & -- downgraded \\
3 & SHMT1/2, MTHFD1L & serine/1C metabolism & +0.001--0.008 & \textbf{DOWN in dysf (all p<0.02)} & weakly essential ($-$0.10 each) & -- downgraded \\
X & MTHFD1 & serine/1C metabolism & +0.008 & \textbf{DOWN in dysf (p<0.02)} & \textbf{pan-essential ($-$0.53)} & -- not a target \\
X & ERBB2 & gastric oncogene & +0.069 & DOWN trend & moderately essential ($-$0.36) & FDA approved (HER2+ GC), but a confound \\
X & RAC1 & NK protrusion GTPase & +0.077 & \textbf{DOWN in dysf (p=0.04)} & \textbf{pan-essential ($-$0.74)} & -- not a target \\
\bottomrule
\end{tabular}}
\end{table*}

\begin{figure*}[htbp]
\centering
\includegraphics[width=\textwidth]{figures/fig3_targets.pdf}
\caption{(a) Tumor-specificity vs NK-dysfunction-correlation scatter. (b) Top-15 tumor-intrinsic candidates. (c) Category composition of the 37-gene pool. \textit{Alt text: Three-panel target-prioritization figure.}}\label{fig:targets}
\end{figure*}

\setcounter{subsection}{6}
\subsection{The graph as a convergent probe of transcriptional reach}

Ablating the \url{metabolic_crosstalk} edge from the enriched heterogeneous graph abolishes the corresponding embedding coupling (H1 embedding-coupling 0.128$\to-$0.002; Table~\ref{tab:ablation}) while leaving the within-axis \url{sm_topology_axis} coupling comparatively intact -- an independent, architecture-based confirmation that the upstream metabolic coupling is transcriptionally absent, converging with the null H1/H2 correlations (\S{}3.2). In the label-masked multi-view audit, learned view-fusion is significantly \textit{worse} than a no-graph baseline for both AUROC and AUPRC ($\Delta=-0.010/-0.021$, bootstrap CI excluding 0), and no view is stably top-ranked across seeds -- the unmasked sensitivity analysis reaches much higher accuracy (0.971/0.936), confirming that label overlap, not graph structure, drove the earlier optimistic result. A node-label permutation test confirms both mechanism-grounded edge types impose their intended geometry (both empirical $P=0.001$), verifying structural specificity without implying biological or predictive validation. Together, three methodologically independent lines of evidence -- correlation, graph ablation, and multi-view audit -- converge on the same boundary: the effector layer is transcriptionally present; the upstream metabolic coupling is not. Full multi-view diagnostics are in Supplementary Results S3.7 and Fig.~S1.

\begin{table*}[htbp]
\centering
\caption{Graph ablation on the enriched real graph. Embedding-coupling = mean cosine similarity between module gene embeddings; modularity = intra-module minus inter-module cosine similarity.}\label{tab:ablation}
\footnotesize
\begin{tabularx}{\textwidth}{@{}lXXXX@{}}
\toprule
Variant & Edges & Modularity & H1 embedding-coupling (serine$\leftrightarrow$SM) & H2 embedding-coupling (SM$\leftrightarrow$protrusion) \\
\midrule
FULL (all edges) & 1,612 & 0.279 & 0.128 & 0.183 \\
$-$MC (without \url{metabolic_crosstalk}) & 1,312 & 0.288 & $-$0.002 & 0.146 \\
$-$SST (mechanism edges removed; generic priors only) & 492 & 0.287 & $-$0.002 & $-$0.000 \\
\bottomrule
\end{tabularx}
\end{table*}

Beyond the two mechanism cards run end-to-end above, an exploratory scan applying the same comparisons to all four registered cards (including the two validated only structurally) across additional cohorts and modalities is reported in Supplementary Results S3.8 and Fig.~S2; its pre-registered cross-mechanism gate is not met (\url{comparative_atlas_only}), so it is not treated as further evidence for or against the reach boundary above -- only the two-card result in \S{}3.7/\S{}4.1 is.

\section{Discussion}

\subsection{Summary of findings}

Our central result is a \textbf{map of the mechanism's transcriptional reach}, not a blanket claim of recovery. Four findings anchor it:

\begin{enumerate}
\item \textbf{The effector layer is transcriptionally recoverable in bulk transcriptomes and generalizes to gastric cancer}, but is substantially NK-abundance-driven: the protrusion-machinery$\to$cytotoxicity coupling replicates from independent liver (r=0.55) and gastric transcriptomes, yet after adjusting for NK-fraction it attenuates by $\sim$50\%, holding in two of three bulk cohorts and not the third. The single-cell version of the same coupling does not survive residualization against library size/latent structure or a random-module permutation baseline (\S{}3.2), so we anchor this finding on the bulk result.
\item \textbf{The metabolic coupling is not transcriptionally detectable -- by two independent methods.} Correlation analysis (H1, H2 both null after correction) and graph ablation (removing the \url{metabolic_crosstalk} edge abolishes the imposed embedding coupling, \S{}3.7) converge: this coupling is real biology operating at the metabolite level, not the transcript level, exactly the boundary the mechanism card's ``gated physical ground-truth'' section is designed to mark.
\item \textbf{Transcription does not substitute for the physical topology phenotype -- a finding, not a failure.} Machinery-gene transcription runs opposite to the physical protrusion collapse in intratumoral NK (higher transcript, lower function), reproducing why the anchor lab required single-cell mass spectrometry and super-resolution imaging, and identifying the post-transcriptional membrane-recruitment layer as the critical gap for future experimental work.
\item \textbf{The transcriptional-reach boundary generalizes across mechanisms.} A second, independently authored mechanism card (TGF$\beta$$\to$SMAD$\to$NK exclusion [10]) run end-to-end on the same cohorts reproduces the pattern: NK-effector/outcome co-variation is recoverable, but the mechanism-specific causal predictions (TGF$\beta$ signaling suppresses activating receptors; CAF-ECM excludes cytotoxic NK) fail in bulk, once NK-cell fraction is controlled -- the same shared cause as the Zheng card's upstream metabolic coupling failing.
\end{enumerate}

Independent orthogonal validation (NK-state DE + DepMap essentiality, \S{}3.5) refines the 37-gene candidate list into evidence tiers: SMPD3 and SMPD1 pass both filters cleanly; SGMS2 and NT5E/CD73 pass with caveats (discordant external DE and label self-referentiality, respectively); five serine-pathway genes and ERBB2/RAC1 are downgraded or excluded. A per-sample readout combining NK SM-catabolism score and HAVCR2 (Tim3) expression can stratify samples by the SM-restoration + Tim3-blockade combination logic proposed in the anchor paper -- an in-silico hypothesis for experimental testing, not a validated clinical predictor; this is consistent with an independent HCC study reporting co-regulation of NK checkpoints (HAVCR2, TIGIT, PDCD1, LAG-3, CLEC12B) [46].

\subsection{The mechanism-card approach}

A key contribution is the mechanism-card abstraction itself: a machine-readable YAML formalism separating a mechanism's computational operationalization from the pipeline that executes it, so applying the framework to a new mechanism requires only authoring a new card, not rewriting the pipeline core. Beyond reuse, the card format enforces scientific discipline shown here to be load-bearing rather than decorative: explicit claim boundaries (letting us state without ambiguity that machinery transcription does not proxy physical topology), gated physical ground-truth targets (preventing conflation of transcriptional proxy with direct measurement), and pre-registered hypotheses with a defined recovery criterion (so Arm A's partial outcome is reported as a structured result rather than reframed).

The mechanism-card registry holds four cards (serine--SM, adenosine-mediated suppression, TGF$\beta$-driven exclusion, MICA/B shedding) sharing one YAML schema and low pairwise gene overlap (Jaccard 0.07--0.19). \textbf{Two have now been run end-to-end on real data}: the TGF$\beta$$\to$SMAD$\to$NK-exclusion card reproduces the same transcriptional-reach boundary as the Zheng card rather than a card-specific artifact -- effector/outcome co-variation recovers, but its mechanism-specific causal predictions fail in bulk. The adenosine and MICA/B cards remain structurally validated but not yet run end-to-end. Reusability is therefore an established capability, not only a design aspiration.

\subsection{Limitations}

The complete 17-item limitations discussion is in Supplementary Results S4.3; key points are summarized here.

\begin{enumerate}
\item \textbf{Transcriptional proxy $\neq$ physical topology}, and serine/SM crosstalk is a metabolite-level effect not captured by transcription -- both now empirically supported by the H1/H2 null results.
\item \textbf{No experimental validation}; all targets are computationally prioritized, with recommended assays bridging to a companion wet-lab program.
\item \textbf{Single-cell pseudoreplication and technical confounding.} The H3 single-cell coupling requires count-depth/module-permutation controls beyond pseudoreplication correction alone (\S{}3.2), and does not survive them; a leave-one-sample-out check shows the pooled estimate itself is not an artifact of one sample.
\item \textbf{The graph does not outperform top tabular or signature baselines on accuracy} -- confirmed by two independent stricter audits (\S{}3.4, \S{}3.7); an exploratory single-edge cross-cohort transfer test proved highly seed-sensitive and supports no directional claim.
\item \textbf{Residual NK-side bias in the tumor-intrinsic candidate pool} (46\% retain NK-side module annotation) and \textbf{omission of intracellular/TF-level regulators} (e.g. CREM [27]).
\item \textbf{No clinical-outcome anchor}; NK states are linked to scRNA-defined labels, not survival or therapy response.
\item \textbf{Multi-card validation covers two of four cards end-to-end}; the bulk effector coupling is substantially NK-abundance-driven (\S{}3.2); DepMap essentiality uses the real, current 26Q1 release; NK-state DE is underpowered for the dysfunctional group (n=20) though 88\% externally concordant.
\item \textbf{The NK-state classification label partly overlaps the classifier's own input features} (a gene-set separation audit found label-defining marker genes are also present in the input expression vector), so \S{}3.4's classification accuracy should be read accordingly; the SST-axis mechanism tests (H1--H5) do not share this issue.
\item \textbf{No hard per-cell QC threshold was applied} to the scRNA-seq data; technical variance is instead addressed via count-depth residualization.
\end{enumerate}

\subsection{Future directions}

\begin{itemize}
\item \textbf{Additional cancer types and mechanisms.} Other digestive-tract cancers (colorectal, pancreatic, esophageal) using the same card; the registry already carries authored-but-not-run cards for adenosine suppression, MICA/B-ADAM17 shedding, two checkpoint axes [47,48], and two metabolite axes, plus a not-yet-authored CLEC12B-lipoprotein-lipase axis [46].
\item \textbf{Physical topology integration} once membrane protrusion/microvilli imaging data become available, and a \textbf{prospective validation cohort} with paired bulk RNA-seq, scRNA-seq, and functional NK assays.
\item \textbf{Genome-scale learnable encoder} (per-node GCN/Transformer as in TREE [49] or GRAFT [50]), reserved for a genome-scale panel and an independently powered endpoint free of the small-panel overfitting risk exposed here.
\end{itemize}

\section{Conclusion}

GC-NKGraph-Atlas demonstrates \textit{how far} a specific, published immune-evasion mechanism -- the serine--sphingomyelin--membrane-topology axis of NK dysfunction -- can be surveyed at cohort scale from public transcriptomes. Using a single-cell-informed heterogeneous graph framework and a two-arm design, we establish a three-layer scoping map: (i) the effector arm is recoverable in bulk transcriptomes from liver to gastric cancer, but $\sim$50\% of the bulk coupling is NK-cell-abundance covariation and the single-cell version does not survive technical-confound controls, so we report it as partially recoverable and substantially infiltration-driven rather than robustly recovered; this boundary reproduces on a second, independently authored mechanism card (TGF$\beta$$\to$SMAD$\to$NK exclusion); (ii) the upstream metabolic coupling is \textbf{not} transcriptionally detectable after correction ($P$=0.20), consistent with metabolite-level regulation -- independently confirmed by graph ablation; and (iii) the physical topology phenotype is fundamentally disconnected from machinery transcription -- a finding, not a failure, that benchmarks the natural limit of transcriptome-based reconstruction for membrane-lipid mechanisms. This map, not an overclaim of full reconstruction, is the contribution.

The mechanism-card abstraction, demonstrated here on the Zheng 2023 serine--SM axis, provides a disciplined formalism for converting a published mechanism into a scalable study with explicit claim boundaries and gated physical ground-truth. This discipline is load-bearing: it is precisely the card's separation of computable-from-physical that lets us report the topology disconnection as a structured result, and the pre-registered recovery definition that forces the H2 non-significance to be reported plainly. The framework has now been run end-to-end on two independently authored cards, which converge on the same transcriptional-reach boundary; extending it to the remaining registered cards (adenosine-mediated NK suppression, MICA/B shedding) is a natural next step, so multi-card reuse is an established capability to be broadened rather than only a design aspiration.

\section*{Data Availability}

All code and configuration files are available at \url{https://github.com/nblvguohao/GC-NKGraph-Atlas}. The mechanism-card template and the Zheng 2023 NK SM-topology card are provided under \url{configs/mechanism_cards/}. Synthetic test data can be generated via python \url{src/common/synthetic_data.py}.

\textbf{Reproducibility.} The pipeline supports a fully self-contained synthetic data mode ($\approx$17 MB, no real patient data) that exercises all phases end-to-end:
\begin{verbatim}
pip install -r requirements.txt
python src/common/synthetic_data.py  # -> data/synthetic/
python src/pipeline.py --synthetic --force
\end{verbatim}
An environment lock file (\url{environment.yml}) pins all dependency versions (testable on Linux/macOS/Windows with conda). Unit tests covering graph construction, NK scoring, SST-axis computation, and target prioritization are provided under \url{tests/} and can be run with \texttt{pytest tests/}. The pseudoreplication correction analysis is implemented in \url{src/topology/pseudoreplication_correction.py} and the count-depth control in \url{src/topology/count_depth_control.py}.

\textbf{Real data availability:}

\begin{itemize}
\item TCGA-STAD and TCGA-LIHC: available from the Genomic Data Commons (\url{https://portal.gdc.cancer.gov/})
\item GSE62254, GSE84437, GSE246662: available from the Gene Expression Omnibus (\url{https://www.ncbi.nlm.nih.gov/geo/})
\item DepMap Public 26Q1 (\url{CRISPRGeneEffect.csv}, \url{Model.csv}): the real, current release used for Table~\ref{tab:targets}'s CERES values, obtained directly from the DepMap portal at \url{https://depmap.org/portal/} (interactive download required -- see Limitations). An earlier pass of this analysis used DepMap Public 24Q2 via the public figshare API (\url{https://plus.figshare.com/articles/dataset/DepMap_24Q2_Public/25880521}) before 26Q1 access was available; that snapshot has been superseded throughout.
\item No novel sequencing data were generated for this study.
\end{itemize}

\textbf{Pre-registration:} The hypotheses H1--H5 and the original full-recovery criterion for the positive-control arm were registered in \url{configs/sst_axis_config.yaml} before execution. Because the full criterion was not met, the manuscript reports the result as a partial-recovery scoping map and explicitly separates pre-registered outcomes from post-hoc interpretation.

\section*{Author Contributions}

\textbf{Guohao Lyu:} Conceptualization, Methodology, Software, Formal analysis, Investigation, Data curation, Visualization, Writing -- original draft. \textbf{Yingchun Xia:} Methodology, Software, Validation, Writing -- review \& editing. \textbf{Huichao Liu:} Software, Data curation, Validation. \textbf{Xiaolei Zhu:} Formal analysis, Methodology, Writing -- review \& editing. \textbf{Shuai Yang:} Investigation, Validation, Visualization. \textbf{Ailian Zhou:} Conceptualization, Supervision, Funding acquisition, Writing -- review \& editing. \textbf{Lichuan Gu:} Conceptualization, Supervision, Project administration, Funding acquisition, Writing -- review \& editing.

All authors read and approved the final manuscript.

\section*{Ethics Approval and Consent to Participate}

Not applicable. This study used only publicly available, de-identified datasets (TCGA-LIHC, TCGA-STAD, GSE62254, GSE84437, GSE246662) obtained from the Genomic Data Commons and the Gene Expression Omnibus. No new human or animal subjects were involved, and no identifiable personal data were generated or analyzed. Use of these public data complies with the respective data-access policies.

\section*{Competing Interests}

The authors declare that they have no competing interests.

\section*{Funding}

This work was supported by grants from the National Natural Science Foundation of China (32472007, 62301006, 62301008), the Natural Science Foundation of Anhui Province (2308085MF217, 2308085QF202), and the Anhui Province Key Laboratory of Intelligent Agricultural Technology and Equipment.


\begin{thebibliography}{99}
\bibitem[1]{ref1} Vivier E, Tomasello E, Baratin M, et al. Functions of natural killer cells. \textit{Nat Immunol} 2008;9:503--10.
\bibitem[2]{ref2} Chiossone L, Dumas PY, Vienne M, Vivier E. Natural killer cells and other innate lymphoid cells in cancer. \textit{Nat Rev Immunol} 2018;18:671--88.
\bibitem[3]{ref3} Huntington ND, Cursons J, Rautela J. The cancer--natural killer cell immunity cycle. \textit{Nat Rev Cancer} 2020;20:437--54.
\bibitem[4]{ref4} Myers JA, Miller JS. Exploring the NK cell platform for cancer immunotherapy. \textit{Nat Rev Clin Oncol} 2021;18:85--100.
\bibitem[5]{ref5} Shimasaki N, Jain A, Campana D. NK cells for cancer immunotherapy. \textit{Nat Rev Drug Discov} 2020;19:200--18.
\bibitem[6]{ref6} O'Brien KL, Finlay DK. Immunometabolism and natural killer cell responses. \textit{Nat Rev Immunol} 2019;19:282--90.
\bibitem[7]{ref7} Terr\'en I, Orrantia A, Vitall\'e J, et al. NK cell metabolism and tumor microenvironment. \textit{Front Immunol} 2019;10:2278.
\bibitem[8]{ref8} Laskowski TJ, Biederst\"adt A, Rezvani K. Natural killer cells in antitumour adoptive cell immunotherapy. \textit{Nat Rev Cancer} 2022;22:557--75.
\bibitem[9]{ref9} Andr\'e P, Denis C, Soulas C, et al. Anti-NKG2A mAb is a checkpoint inhibitor that promotes anti-tumor immunity by unleashing both T and NK cells. \textit{Cell} 2018;175:1731--43.e13.
\bibitem[10]{ref10} Gao J, Zheng X, Liu Y, et al. TGF-$\beta$ impairs NK cell migration and cytotoxicity. \textit{J Immunother Cancer} 2023;11:e005785.
\bibitem[11]{ref11} Melaiu O, Lucarini V, Cifaldi L, Fruci D. Influence of the tumor microenvironment on NK cell function in solid tumors. \textit{Front Immunol} 2020;10:3038.
\bibitem[12]{ref12} Zheng X, Hou Z, Qian Y, et al. Tumors evade immune cytotoxicity by altering the surface topology of NK cells. \textit{Nat Immunol} 2023;24:748--59. doi:10.1038/s41590-023-01462-9.
\bibitem[13]{ref13} Possemato R, Marks KM, Shaul YD, et al. Functional genomics reveal that the serine synthesis pathway is essential in breast cancer. \textit{Nature} 2011;476:346--50.
\bibitem[14]{ref14} Locasale JW. Serine, glycine and one-carbon units: cancer metabolism in full circle. \textit{Nat Rev Cancer} 2013;13:572--83.
\bibitem[15]{ref15} Ogretmen B. Sphingolipid metabolism in cancer -- signalling and drug resistance. \textit{Nat Rev Cancer} 2018;18:33--50.
\bibitem[16]{ref16} Hannun YA, Obeid LM. Sphingolipids and their metabolism in physiology and disease. \textit{Nat Rev Mol Cell Biol} 2018;19:175--91.
\bibitem[17]{ref17} Zheng X, Tian Z, Wei H. Sphingomyelin is a prospective metabolic immune checkpoint for natural killer cells. \textit{Clin Transl Med} 2023;13:e1395.
\bibitem[18]{ref18} Newman AM, Steen CB, Liu CL, et al. Determining cell type abundance and expression from bulk tissues with digital cytometry. \textit{Nat Biotechnol} 2019;37:773--82.
\bibitem[19]{ref19} Finotello F, Mayer C, Plattner C, et al. Molecular and pharmacological modulators of the tumor immune contexture revealed by deconvolution of RNA-seq data. \textit{Genome Med} 2019;11:34.
\bibitem[20]{ref20} Sun D, Guan X, Moran AE, et al. Identifying phenotype-associated subpopulations by integrating bulk and single-cell sequencing data. \textit{Nat Biotechnol} 2022;40:527--38.
\bibitem[21]{ref21} Bai Z, Feng B, McClory SE, et al. Single-cell CAR T atlas reveals type 2 function in 8-year leukaemia remission. \textit{Nature} 2024;634:702--11. doi:10.1038/s41586-024-07762-w.
\bibitem[22]{ref22} Jin S, Guerrero-Juarez CF, Zhang L, et al. Inference and analysis of cell--cell communication using CellChat. \textit{Nat Commun} 2021;12:1088.
\bibitem[23]{ref23} Efremova M, Vento-Tormo M, Teichmann SA, Vento-Tormo R. CellPhoneDB: inferring cell--cell communication from combined expression of multi-subunit ligand--receptor complexes. \textit{Nat Protoc} 2020;15:1484--506.
\bibitem[24]{ref24} Browaeys R, Saelens W, Saeys Y. NicheNet: modeling intercellular communication by linking ligands to target genes. \textit{Nat Methods} 2020;17:159--62.
\bibitem[25]{ref25} Alghamdi N, Chang W, Dang P, et al. A graph neural network model to estimate cell-wise metabolic flux using single-cell RNA-seq data. \textit{Genome Res} 2021;31:1867--84.
\bibitem[26]{ref26} Gulati GS, Sikandar SS, Wesche DJ, et al. Single-cell transcriptional diversity is a hallmark of developmental potential. \textit{Science} 2020;367:405--11.
\bibitem[27]{ref27} Rafei H, Basar R, Acharya S, et al. CREM is a regulatory checkpoint of CAR and IL-15 signalling in NK cells. \textit{Nature} 2025;643:1076--86. doi:10.1038/s41586-025-09087-8.
\bibitem[28]{ref28} Kipf TN, Welling M. Semi-supervised classification with graph convolutional networks. \textit{Proc ICLR} 2017.
\bibitem[29]{ref29} Veli\v{c}kovi\'c P, Cucurull G, Casanova A, et al. Graph attention networks. \textit{Proc ICLR} 2018.
\bibitem[30]{ref30} Schlichtkrull M, Kipf TN, Bloem P, et al. Modeling relational data with graph convolutional networks. \textit{Proc ESWC} 2018:593--607.
\bibitem[31]{ref31} Wang T, Shao W, Huang Z, et al. MOGONET integrates multi-omics data using graph convolutional networks allowing patient classification and biomarker identification. \textit{Nat Commun} 2021;12:3445.
\bibitem[32]{ref32} Hu Z, Dong Y, Wang K, Sun Y. Heterogeneous graph transformer. \textit{Proc Web Conf} 2020:2704--10.
\bibitem[33]{ref33} Stuart T, Butler A, Hoffman P, et al. Comprehensive integration of single-cell data. \textit{Cell} 2019;177:1888--902.
\bibitem[34]{ref34} Hao Y, Hao S, Andersen-Nissen E, et al. Integrated analysis of multimodal single-cell data. \textit{Cell} 2021;184:3573--87.e29.
\bibitem[35]{ref35} Gayoso A, Lopez R, Xing G, et al. A Python library for probabilistic analysis of single-cell omics data. \textit{Nat Biotechnol} 2022;40:163--66.
\bibitem[36]{ref36} Wolf FA, Angerer P, Theis FJ. SCANPY: large-scale single-cell gene expression data analysis. \textit{Genome Biol} 2018;19:15.
\bibitem[37]{ref37} Luecken MD, Theis FJ. Current best practices in single-cell RNA-seq analysis: a tutorial. \textit{Mol Syst Biol} 2019;15:e8746.
\bibitem[38]{ref38} Lotfollahi M, Naghipourfar M, Luecken MD, et al. Mapping single-cell data to reference atlases by transfer learning. \textit{Nat Biotechnol} 2022;40:121--30.
\bibitem[39]{ref39} Subramanian A, Tamayo P, Mootha VK, et al. Gene set enrichment analysis: a knowledge-based approach for interpreting genome-wide expression profiles. \textit{Proc Natl Acad Sci USA} 2005;102:15545--50.
\bibitem[40]{ref40} Liberzon A, Birger C, Thorvaldsd\'ottir H, et al. The Molecular Signatures Database (MSigDB) hallmark gene set collection. \textit{Cell Syst} 2015;1:417--25.
\bibitem[41]{ref41} Fey M, Lenssen JE. Fast graph representation learning with PyTorch Geometric. \textit{ICLR Workshop} 2019.
\bibitem[42]{ref42} Chen T, Guestrin C. XGBoost: a scalable tree boosting system. \textit{Proc KDD} 2016:785--94.
\bibitem[43]{ref43} Ke G, Meng Q, Finley T, et al. LightGBM: a highly efficient gradient boosting decision tree. \textit{Adv NeurIPS} 2017:3146--54.
\bibitem[44]{ref44} Cancer Genome Atlas Research Network. Comprehensive molecular characterization of gastric adenocarcinoma. \textit{Nature} 2014;513:202--9.
\bibitem[45]{ref45} Cristescu R, Lee J, Nebozhyn M, et al. Molecular analysis of gastric cancer identifies subtypes associated with distinct clinical outcomes. \textit{Nat Med} 2015;21:449--56.
\bibitem[46]{ref46} Sun P, Xu X, Hu B, et al. Targeting NK cell CLEC12B enhances cancer immunotherapy. \textit{Nat Immunol} 2026;27:985--99. doi:10.1038/s41590-026-02471-0.
\bibitem[47]{ref47} Barry KC, Hsu J, Broz ML, et al. A natural killer--dendritic cell axis defines checkpoint therapy--responsive tumor microenvironments. \textit{Nat Med} 2018;24:1178--91.
\bibitem[48]{ref48} Zhang Q, Bi J, Zheng X, et al. Blockade of the checkpoint receptor TIGIT prevents NK cell exhaustion and elicits potent anti-tumor immunity. \textit{Nat Immunol} 2018;19:723--32.
\bibitem[49]{ref49} Su X, Hu P, Li D, et al. Interpretable identification of cancer genes across biological networks via transformer-powered graph representation learning. \textit{Nat Biomed Eng} 2025;9:371--89.
\bibitem[50]{ref50} Cho SP, Cho YR. GRAFT: a graph-aware fusion transformer for cancer driver gene prediction. \textit{Brief Bioinform} 2026;27:bbaf706.
\end{thebibliography}

\end{document}
